\chapter{Le desk de trading côté buy side et ses courtiers}\label{ch:mx:the-buy-side-trading-desk-and-its-brokers}

Le desk d'un fonds envoie ses ordres aux algorithmes de six courtiers et les rémunère en commissions. Lequel des courtiers est le meilleur ne se voit pas dans les données d'un trimestre~: le desk confie donc le choix à une
\omterm{def:mx:the-buy-side-trading-desk-and-its-brokers:wheel}{algo wheel} et mesure. Ce chapitre décrit ce que fait le desk et les systèmes avec lesquels il travaille, puis simule plus d'un an de ses flux pour
répondre à trois questions~: combien de mois d'ordres faut-il pour distinguer le meilleur courtier du second, ce que gagne et perd une
\omterm{def:mx:the-buy-side-trading-desk-and-its-brokers:wheel}{algo wheel} pondérée par la performance, et ce qu'il advient du classement des courtiers lorsque les traders choisissent qui reçoit les ordres difficiles.

\section{Ce que fait le desk}

Un gérant de portefeuille décide ce qu'il faut acheter et vendre~; le desk décide de la manière. Il reçoit les ordres, choisit pour chacun un courtier, un algorithme et
ses paramètres (ou une négociation de bloc), les suit pendant leur exécution et répond du résultat au titre de l'obligation de meilleure exécution qui pèse sur la société
(livre~1, chapitre~11). Deux systèmes portent ce travail.

\begin{definition}[Système de gestion des ordres, système de gestion de l'exécution]\label{def:mx:the-buy-side-trading-desk-and-its-brokers:oms}
Un \emph{système de gestion des ordres}\index{système de gestion des ordres} (OMS) tient les ordres et les positions du fonds, de la décision du gérant de portefeuille jusqu'au
règlement-livraison~: contrôles de conformité, affectations aux comptes, historique de ce qui a été commandé et exécuté. Un
\emph{système de gestion de l'exécution}\index{système de gestion de l'exécution} (EMS) est l'écran du trader pour suivre les ordres en cours d'exécution~: données de marché, algorithmes des courtiers, accès direct au marché (DMA), ainsi que les \omterm{def:mx:the-almgren-chriss-framework:parent}{ordres enfants} et les
exécutions en temps réel.
\end{definition}

L'OMS connaît l'instant de décision et le prix~; l'EMS connaît chaque \omterm{def:mx:the-almgren-chriss-framework:parent}{ordre enfant}. Le TCA (chapitre~19) a besoin des deux, et le défaut le plus courant dans les données d'un desk est que
l'instant de décision n'a jamais quitté la tête du gérant de portefeuille.

\section{Trading high-touch et trading low-touch}

\begin{definition}[Trading high-touch et trading low-touch]\label{def:mx:the-buy-side-trading-desk-and-its-brokers:touch}
Le \emph{trading high-touch}\index{trading high-touch} confie un ordre à une personne chez le courtier (un sales-trader ou un cash trader, livre~1, chapitre~2),
qui le traite avec discernement, trouve des contreparties naturelles pour les blocs et peut engager le capital du courtier. \emph{Trading
low-touch}\index{trading low-touch} l'envoie aux canaux électroniques du courtier (algorithmes, accès direct au marché, routage intelligent), le trader
buy side en faisant le choix et en assurant la supervision.
\end{definition}

Le \omterm{def:mx:the-buy-side-trading-desk-and-its-brokers:touch}{trading high-touch} convient aux ordres dont la taille ou l'information font de la réaction du marché le principal coût~: un gros bloc sur un titre peu liquide, une transaction que le fonds
ne veut pas exposer, à l'aide des \omterm{def:mx:dark-pools-and-blocks:ioi}{indications d'intérêt} (IOI) et des \omterm{def:mx:dark-pools-and-blocks:conditional}{ordres conditionnels} du chapitre~9. Le \omterm{def:mx:the-buy-side-trading-desk-and-its-brokers:touch}{trading low-touch} concerne le reste, soit la majorité des ordres en nombre. La frontière
se déplace avec la difficulté de l'ordre, et la comparaison des courtiers aussi~: un courtier qui reçoit les ordres difficiles paraît cher, quoi qu'il fasse.

\section{Les algo wheels}

\begin{definition}[Algo wheel]\label{def:mx:the-buy-side-trading-desk-and-its-brokers:wheel}
Une \emph{algo wheel}\index{algo wheel} répartit les ordres low-touch du desk entre les algorithmes des courtiers selon une règle plutôt que selon le choix du trader~:
au hasard, à tour de rôle, en stratifiant selon la difficulté de l'ordre ou en pondérant par la performance passée, de sorte que les résultats de chaque courtier puissent être comparés sur
un flux comparable.
\end{definition}

Une \omterm{def:mx:the-buy-side-trading-desk-and-its-brokers:wheel}{algo wheel} est un test A/B (livre~7, chapitre~21) dont l'unité de randomisation est l'ordre. Le desk simulé de ce chapitre envoie 600 ordres par
mois à six courtiers. La difficulté de chaque ordre est son \omterm{def:mx:transaction-cost-analysis:pretrade}{estimation des coûts pré-négociation} (chapitre~19), suivant une loi lognormale de moyenne 15 points de base et d'écart type
de 25~; son coût est l'effet du courtier, plus cette difficulté, plus un bruit suivant une loi de Student $t$ à quatre degrés de liberté et d'écart type
de 30 points de base (les coûts d'exécution réels sont à queue épaisse). Les effets des courtiers valent 0, 3, 4, 5, 6 et 8 points de base~: le courtier~1 est le meilleur, le courtier~2 est moins bon de 3
points de base.

L'effet minimal détectable (livre~7, chapitre~21) fixe la durée pendant laquelle l'\omterm{def:mx:the-buy-side-trading-desk-and-its-brokers:wheel}{algo wheel} doit tourner. Pour détecter une différence $\delta$ entre deux courtiers avec
un test bilatéral au seuil de 5~\% et une puissance de 80~\% (livre~4, chapitre~12), chacun requiert $2\sigma^2(z_{0.975}+z_{0.8})^2/\delta^2$ ordres~; avec six
courtiers se partageant 600 ordres par mois, $\delta=3$ et $\sigma=30$, soit 15.7 mois. Sans ajustement sur la difficulté, le bruit inclut celle des ordres
$\sigma=\sqrt{30^2+25^2}=39.1$, et il faut 26.6 mois.

\omcode{code/firm/algowheel/firm_algowheel.py}{78}{82}{Le nombre de mois nécessaires à une algo wheel uniforme pour détecter une différence entre deux courtiers~: les ordres requis par courtier pour un test bilatéral à la puissance choisie, multipliés par le nombre de courtiers, rapportés au flux mensuel du desk.}\label{lst:mx:the-buy-side-trading-desk-and-its-brokers:months}

La \Cref{fig:mx:the-buy-side-trading-desk-and-its-brokers:power} vérifie la formule par simulation~: 300 desks, chacun testant le courtier~2 contre le courtier~1
tous les deux mois. Le test ajusté atteint une puissance de 80~\% après 16 mois, le test brut après 24. Stratifier l'\omterm{def:mx:the-buy-side-trading-desk-and-its-brokers:wheel}{algo wheel} selon la difficulté (chaque courtier reçoit le
même mélange d'ordres dans vingt classes de difficulté) rend la comparaison brute un peu plus précise, avec un écart type de 1.51 points de base après un
an contre 1.59 au hasard, mais l'ajustement fait bien davantage~: il retire la difficulté du bruit au lieu de l'équilibrer.

\begin{omfigure}
\begin{tikzpicture}
\begin{axis}[width=9.6cm,height=5.6cm,xlabel={\small mois de flux (600 ordres par mois)},ylabel={\small puissance (\%)},
  tick label style={font=\footnotesize},grid=major,grid style={gray!25},xmin=0,xmax=41,ymin=0,ymax=100,
  legend columns=2,legend style={font=\footnotesize,at={(0.5,-0.3)},anchor=north,draw=none,fill=none,
  /tikz/every even column/.append style={column sep=8pt}},table/col sep=comma]
\addplot[omDef,very thick,mark=*,mark size=1.4pt] table[x=months,y=adjusted]{figdata/microstructure/20-the-buy-side-trading-desk-and-its-brokers/power.csv};
\addplot[omProp,very thick,mark=square*,mark size=1.4pt] table[x=months,y=raw]{figdata/microstructure/20-the-buy-side-trading-desk-and-its-brokers/power.csv};
\addplot[gray,dashed,domain=0:41] {80};
\legend{ajusté pour la difficulté,brut}
\end{axis}
\end{tikzpicture}

\omcaption{La probabilité qu'une \omterm{def:mx:the-buy-side-trading-desk-and-its-brokers:wheel}{algo wheel} uniforme fasse apparaître que le courtier~2 est significativement moins bon que le courtier~1 (une vraie différence de 3 points de base), en fonction des
mois de flux~: 300 desks simulés. Données~: \texttt{mx\_wheel.power\_study}.}\label{fig:mx:the-buy-side-trading-desk-and-its-brokers:power}
\end{omfigure}

Une \omterm{def:mx:the-buy-side-trading-desk-and-its-brokers:wheel}{algo wheel} uniforme paie son savoir~: elle envoie cinq sixièmes du flux à des courtiers qui ne sont pas les meilleurs, soit 4.33 points de base de plus par ordre que si elle
utilisait toujours le courtier~1. Une \omterm{def:mx:the-buy-side-trading-desk-and-its-brokers:wheel}{algo wheel} pondérée par la performance dépense moins. L'échantillonnage de Thompson tient à jour une loi a posteriori du coût ajusté de chaque courtier et envoie chaque ordre au
courtier dont le coût tiré au hasard est le plus bas~; il apprend au fil de l'eau (la \cref{fig:mx:the-buy-side-trading-desk-and-its-brokers:thompson}).
Sur 24 mois, elle a envoyé 87.2~\% des ordres au courtier~1 (97.5~\% le dernier mois), pour un surcoût de 0.54 point de base par ordre. Le prix à payer est l'information~: le courtier~2
a reçu 5.4~\% du flux, et le test du courtier~1 contre le courtier~2 avait une puissance de 88~\% après 24 mois, contre 100~\% pour l'\omterm{def:mx:the-buy-side-trading-desk-and-its-brokers:wheel}{algo wheel} uniforme, qui
connaît tout aussi bien les quatre autres courtiers. Une \omterm{def:mx:the-buy-side-trading-desk-and-its-brokers:wheel}{algo wheel} qui doit aussi garder les courtiers engagés et repérer un courtier qui se dégrade a besoin d'un seuil minimal
de parts d'ordres pour chaque courtier.

\begin{omfigure}
\begin{tikzpicture}
\begin{axis}[width=9.6cm,height=5.4cm,xlabel={\small mois},ylabel={\small part des ordres confiés au courtier 1 (\%)},
  tick label style={font=\footnotesize},grid=major,grid style={gray!25},xmin=0,xmax=25,ymin=0,ymax=100,
  legend columns=2,legend style={font=\footnotesize,at={(0.5,-0.3)},anchor=north,draw=none,fill=none,
  /tikz/every even column/.append style={column sep=8pt}},table/col sep=comma]
\addplot[omDef,very thick,mark=*,mark size=1.2pt] table[x=month,y=thompson]{figdata/microstructure/20-the-buy-side-trading-desk-and-its-brokers/thompson.csv};
\addplot[omProp,very thick] table[x=month,y=uniform]{figdata/microstructure/20-the-buy-side-trading-desk-and-its-brokers/thompson.csv};
\legend{échantillonnage de Thompson,uniforme}
\end{axis}
\end{tikzpicture}

\omcaption{Part des ordres de chaque mois envoyée au meilleur courtier par une \omterm{def:mx:the-buy-side-trading-desk-and-its-brokers:wheel}{algo wheel} à échantillonnage de Thompson et par une \omterm{def:mx:the-buy-side-trading-desk-and-its-brokers:wheel}{algo wheel} uniforme, moyennée sur 60 desks simulés.
Données~: \texttt{mx\_wheel.thompson\_study}.}\label{fig:mx:the-buy-side-trading-desk-and-its-brokers:thompson}
\end{omfigure}

\section{Évaluer les courtiers}

\begin{definition}[Grille d'évaluation des courtiers]\label{def:mx:the-buy-side-trading-desk-and-its-brokers:scorecard}
Une \emph{grille d'évaluation des courtiers}\index{grille d'évaluation des courtiers} est le rapport périodique du desk sur chaque courtier~: ordres et montants échangés, coûts comparés aux
indices de référence convenus, bruts et ajustés de la difficulté avec leurs erreurs types, et des mesures plus souples (service, engagement de capital, qualité de la liquidité des
blocs), à partir desquelles le desk fixe l'allocation de la période suivante.
\end{definition}

\begin{table}[ht]
\centering
\small
\begin{tabular}{@{}lrrrrr@{}}
\toprule
courtier & ordres & coût brut & coût ajusté & e.t. & rang \\
\midrule
courtier 1 & 1\,175 & 15.4 & 15.1 & 0.9 & 1 \\
courtier 2 & 1\,246 & 16.7 & 17.7 & 0.9 & 2 \\
courtier 3 & 1\,171 & 18.6 & 18.7 & 0.9 & 3 \\
courtier 4 & 1\,169 & 21.9 & 21.0 & 0.9 & 5 \\
courtier 5 & 1\,207 & 20.5 & 20.6 & 0.8 & 4 \\
courtier 6 & 1\,232 & 22.9 & 22.7 & 0.8 & 6 \\
\bottomrule
\end{tabular}
\caption{Une année de l'\omterm{def:mx:the-buy-side-trading-desk-and-its-brokers:wheel}{algo wheel} uniforme utilisée comme grille d'évaluation~: coûts en points de base à difficulté moyenne, avec des erreurs types robustes. Les effets vrais sont
0, 3, 4, 5, 6 et 8 points de base au-dessus du courtier~1. Données~: \texttt{mx\_wheel.card}.}\label{tab:mx:the-buy-side-trading-desk-and-its-brokers:card}
\end{table}

La grille d'évaluation du \cref{tab:mx:the-buy-side-trading-desk-and-its-brokers:card} classe les courtiers 4 et 5 dans le mauvais ordre~: leur écart vrai est de 1
point de base, et les erreurs types vont de 0.8 à 0.9. Un classement est une estimation ponctuelle~; la grille doit montrer les erreurs. Anand, Irvine, Puckett et
Venkataraman (2012) ont établi que certains courtiers offrent de meilleures exécutions de façon constante dans le temps, donc que des écarts persistants existent, mais les repérer exige
des années de flux d'un même desk ou les flux regroupés de plusieurs desks.

Sans \omterm{def:mx:the-buy-side-trading-desk-and-its-brokers:wheel}{algo wheel}, la comparaison peut être pire qu'un simple bruit. Supposons que les traders du desk envoient le cinquième le plus difficile des ordres au courtier auquel ils font le plus confiance,
courtier~1, et le reste aux cinq autres. Au bout d'une année, la grille brute classe le courtier~1 en dernier, à 35.1 points de base derrière le courtier~2~; corrigée de
la difficulté, elle le place premier, à 3.8 points de base devant le suivant. L'ajustement ne fonctionne ici que parce que le modèle de difficulté est juste~; avec un
modèle faux, seule la randomisation protège la comparaison.

\section{Payer la recherche et l'exécution}

\begin{definition}[Accord de partage de commissions, dissociation de la recherche]\label{def:mx:the-buy-side-trading-desk-and-its-brokers:csa}
Un \emph{accord de partage de commissions}\index{accord de partage de commissions} (commission sharing agreement) permet à un fonds de verser des commissions à un courtier exécutant et d'en orienter une partie vers
d'autres sociétés qui fournissent de la recherche. La \emph{dissociation de la recherche}\index{dissociation de la recherche} (research unbundling) sépare le paiement de la recherche du paiement de
l'exécution, de sorte que chacun soit tarifé et choisi de façon indépendante.
\end{definition}

Les commissions groupées lient le choix du courtier à la recherche voulue par le fonds, ce qui constitue une autre raison pour laquelle l'allocation du desk peut ne pas suivre la qualité de l'exécution
à elle seule. Les règles ont changé à plusieurs reprises.

\begin{dated}{2026-09}{Payer la recherche}\label{dat:mx:the-buy-side-trading-desk-and-its-brokers:research}
MiFID II, appliquée à partir du 3 janvier 2018, obligeait les prestataires de services d'investissement à séparer les frais de recherche des frais d'exécution~: la recherche payée sur les
ressources propres de la société ou sur un compte de paiement de la recherche convenu avec les clients. Au Royaume-Uni, la FCA a de nouveau autorisé les paiements conjoints pour la recherche et
l'exécution à partir du premier août 2024 (policy statement PS24/9), sous conditions. Dans l'Union européenne, la directive (UE) 2024/2811 du 23 octobre 2024 (qui fait partie du
Listing Act) permet aux sociétés de choisir de payer conjointement ou séparément, à condition de le signaler aux clients et d'évaluer la recherche chaque année~; les États membres devaient l'appliquer
à partir du 6 juin 2026. Aux États-Unis, la dispense de non-intervention du personnel de la SEC du 26 octobre 2017 permettait aux broker-dealers de recevoir des paiements de recherche distincts
provenant de sociétés soumises à MiFID II sans être réglementés comme conseillers en investissement~; prolongée en 2019, elle a expiré le 3 juillet 2023.
\end{dated}

Pour le desk, la \omterm{def:mx:the-buy-side-trading-desk-and-its-brokers:csa}{dissociation de la recherche} a une conséquence pratique~: une fois la recherche payée séparément, le budget d'exécution peut suivre la grille d'évaluation.

\section{Tutoriel~: une algo wheel}\label{tut:mx:the-buy-side-trading-desk-and-its-brokers:tut}

\begin{tutorial}
\textbf{Objectif.} Simuler le flux d'un desk sur six courtiers, lancer les \omterm{def:mx:the-buy-side-trading-desk-and-its-brokers:wheel}{algo wheels} uniforme, stratifiée et à échantillonnage de Thompson, évaluer les courtiers avec et sans
l'ajustement pour la difficulté, et calculer le nombre de mois nécessaires pour les classer.
\textbf{État final~:} les \cref{fig:mx:the-buy-side-trading-desk-and-its-brokers:power,fig:mx:the-buy-side-trading-desk-and-its-brokers:thompson},
le \cref{tab:mx:the-buy-side-trading-desk-and-its-brokers:card} et les nombres des sections 3 et 4.
\begin{tutsteps}
\item \textbf{Flux.} \texttt{mx\_wheel.desk(months, seed)}, \texttt{cost\_of(broker, flow)}.
\item \textbf{Allocation.} \texttt{firm\_algowheel.random\_allocation}, \texttt{stratified\_allocation}, \texttt{Thompson}.
\item \textbf{Évaluation.} \texttt{evaluate(cost, broker, difficulty, k, adjust)} sur la régression avec erreurs types groupées de \texttt{firm.tca}, \texttt{scorecard}~;
\texttt{months\_needed}.
\item \textbf{Études.} \texttt{power\_study()}, \texttt{stratified\_study()}, \texttt{thompson\_study()}, \texttt{routing\_study()}, \texttt{card()}~;
tracer avec \texttt{fig\_wheel.py}.
\end{tutsteps}
\textbf{Pour aller plus loin.} Laisser l'effet d'un courtier dériver dans le temps et voir la vitesse à laquelle chaque \omterm{def:mx:the-buy-side-trading-desk-and-its-brokers:wheel}{algo wheel} le détecte~; ajouter un plancher aux parts de l'échantillonnage de Thompson~;
utiliser CUPED (livre~7, chapitre~21) avec l'estimation pré-négociation comme covariable.
\end{tutorial}

\section{Construction~: algo wheel}\label{bld:mx:the-buy-side-trading-desk-and-its-brokers:algowheel}

\begin{build}
\bfield{Rôle} L'allocation et l'évaluation des courtiers par le desk, utilisées par l'\omterm{def:mx:benchmark-algorithms:algo}{algorithme d'exécution} du chapitre~28 pour comparer ses propres versions.
\bfield{Interface} \texttt{random\_allocation(n, k, rng)}, \texttt{stratified\_allocation(difficulty, k, buckets, rng)},
\texttt{Thompson(k, prior\_mean, prior\_sd, noise\_sd)} avec \texttt{choose} et \texttt{update}, \texttt{evaluate(cost, broker, difficulty, k, adjust, clusters)},
\texttt{months\_needed(delta, sd, orders\_per\_month, k, alpha, power)}, \texttt{scorecard(cost, broker, difficulty, k, names, clusters)}.
\bfield{Règles} Coûts en points de base, positifs lorsqu'ils sont payés~; la difficulté est l'estimation pré-négociation, connue avant l'allocation~; les effets des courtiers sont rapportés à la
difficulté moyenne, avec erreurs types groupées.
\bfield{Tests d'acceptation} \texttt{code/firm/algowheel/tests/}~: la stratification équilibre la difficulté mieux que le hasard~; l'évaluation ajustée retrouve
les effets injectés que l'évaluation brute déforme~; la formule des mois~; l'échantillonnage de Thompson se concentre sur le meilleur courtier.
\bfield{Pour aller plus loin} Effets variables dans le temps~; planchers de parts~; mise en commun entre desks.
\end{build}

\begin{omsources}
\item A.~Anand, P.~Irvine, A.~Puckett et K.~Venkataraman, ``Performance of institutional trading desks: an analysis of persistence in trading
costs'', \emph{Review of Financial Studies} 25(2), 2012.
\item FCA, PS24/9, ``Payment optionality for investment research'', 2024.
\item Directive (EU) 2024/2811 amending Directive 2014/65/EU, 23 octobre 2024.
\item SEC, communiqué de presse 2017-200 (26 octobre 2017), et Commissioner M.~T.~Uyeda, statement on the expiration of the staff no-action letter on
MiFID II, 5 juillet 2023.
\end{omsources}

\section{Exercices}

\begin{exercise}[$\star$]\label{exo:mx:the-buy-side-trading-desk-and-its-brokers:1}
Calculer le nombre d'ordres dont chaque courtier a besoin pour détecter un écart de 3 points de base avec une puissance de 80~\% au seuil de 5~\%, pour un bruit de 30 points de base.
\end{exercise}

\begin{exercise}[$\star$]\label{exo:mx:the-buy-side-trading-desk-and-its-brokers:2}
Le desk ne retient plus que trois courtiers. Calculer le nombre de mois dont la comparaison ajustée des deux meilleurs a désormais besoin.
\end{exercise}

\begin{exercise}[$\star$]\label{exo:mx:the-buy-side-trading-desk-and-its-brokers:3}
Calculer le surcoût par ordre d'une \omterm{def:mx:the-buy-side-trading-desk-and-its-brokers:wheel}{algo wheel} uniforme par rapport au choix systématique du courtier~1.
\end{exercise}

\begin{exercise}[$\star\star$]\label{exo:mx:the-buy-side-trading-desk-and-its-brokers:4}
Expliquer pourquoi la stratification n'améliore que très peu la comparaison brute ici.
\end{exercise}

\begin{exercise}[$\star\star$]\label{exo:mx:the-buy-side-trading-desk-and-its-brokers:5}
Le Thompson sampling permet d'économiser 3.8 points de base par ordre. Que sacrifie-t-il, et dans quel cas un desk préférerait-il l'\omterm{def:mx:the-buy-side-trading-desk-and-its-brokers:wheel}{algo wheel} uniforme~?
\end{exercise}

\begin{exercise}[$\star\star$]\label{exo:mx:the-buy-side-trading-desk-and-its-brokers:6}
Pourquoi le classement brut de la grille a-t-il placé le courtier~1 en dernière position alors que les traders lui envoyaient les ordres les plus difficiles, et sur quelle hypothèse repose l'ajustement~?
\end{exercise}

\begin{exercise}[$\star\star\star$]\label{exo:mx:the-buy-side-trading-desk-and-its-brokers:7}
\emph{Programmation.} L'écart type du bruit étant divisé par deux (15 points de base), combien de mois les comparaisons ajustée et brute demandent-elles selon la formule~?
\end{exercise}

\begin{exercise}[$\star\star\star$]\label{exo:mx:the-buy-side-trading-desk-and-its-brokers:8}
\emph{Trouver la faille.} «~Après un trimestre de rotation, le courtier~3 coûte 2 points de base de moins que le courtier~2, donc on relève son flux.~»
\end{exercise}

\section{Problème~: Six courtiers, une rotation}

\begin{problem}[{Problème du week-end --- six courtiers, une rotation}]\label{pb:mx:the-buy-side-trading-desk-and-its-brokers:1}
Le responsable d'un desk demande combien de temps sa nouvelle rotation doit tourner avant qu'il puisse désigner le meilleur courtier, et s'il faut pondérer la rotation par la performance.

\medskip\noindent\textbf{Partie I --- Le desk.}
\begin{enumerate}
\item Que contiennent respectivement l'OMS et l'EMS, et pourquoi la TCA a-t-elle besoin des deux~?
\item Distinguer le \omterm{def:mx:the-buy-side-trading-desk-and-its-brokers:touch}{trading high-touch} et le \omterm{def:mx:the-buy-side-trading-desk-and-its-brokers:touch}{trading low-touch}, et dire quels ordres vont dans chacun.
\item Pourquoi la difficulté des ordres reçus par un courtier fausse-t-elle sa grille~?
\item Définir une \omterm{def:mx:the-buy-side-trading-desk-and-its-brokers:wheel}{algo wheel} et son unité de randomisation.
\end{enumerate}

\medskip\noindent\textbf{Partie II --- Combien de temps.}
\begin{enumerate}[resume]
\item Décrire le flux simulé et les effets propres à chaque courtier.
\item Calculer le nombre d'ordres dont chaque courtier a besoin pour un test bilatéral à 5~\% avec une puissance de 80~\%.
\item Pourquoi la comparaison brute demande-t-elle davantage de mois que la comparaison ajustée~?
\item \emph{Énoncer le résultat nommé}~: les mois de flux nécessaires pour distinguer le meilleur courtier du second avec une puissance de 80~\%, avec et sans ajustement de la difficulté.
adjustment.
\item Qu'apporte la stratification~?
\end{enumerate}

\medskip\noindent\textbf{Partie III --- Pondération par la performance.}
\begin{enumerate}[resume]
\item Comment le Thompson sampling répartit-il le flux~?
\item Donner sa part de flux allouée au meilleur courtier, son surcoût et sa puissance, par comparaison avec ceux de la rotation uniforme.
\item Que ajouteriez-vous pour protéger le desk contre un courtier dont la qualité se dégrade~?
\item Lire la grille~: que peut-on conclure et que ne peut-on pas conclure au sujet des courtiers 4 et 5~?
\end{enumerate}

\medskip\noindent\textbf{Partie IV --- Jugement et argent.}
\begin{enumerate}[resume]
\item Qu'a fait le routage des traders au classement brut, et que l'ajustement a-t-il permis de rétablir~?
\item Que ont trouvé Anand, Irvine, Puckett et Venkataraman au sujet de la persistance~?
\item Définir un \omterm{def:mx:the-buy-side-trading-desk-and-its-brokers:csa}{accord de partage de commissions} et la \omterm{def:mx:the-buy-side-trading-desk-and-its-brokers:csa}{dissociation de la recherche}.
\item Résumer où en sont les règles de paiement de la recherche au Royaume-Uni, dans l'UE et aux États-Unis.
\item Pourquoi la dissociation compte-t-elle pour la rotation~?
\item Que diriez-vous au responsable du desk~?
\item En une phrase~: pourquoi un desk randomise-t-il~?
\end{enumerate}
\end{problem}

\section{Questions d'entretien}

\begin{interviewq}[$\star$ \iqroles{trader}]\label{iq:mx:the-buy-side-trading-desk-and-its-brokers:1}
Quand donneriez-vous un ordre à un sales-trader plutôt qu'à un algorithme~?
\end{interviewq}

\begin{interviewq}[$\star\star$ \iqroles{chercheur}]\label{iq:mx:the-buy-side-trading-desk-and-its-brokers:2}
Comment concevriez-vous l'évaluation d'une \omterm{def:mx:the-buy-side-trading-desk-and-its-brokers:wheel}{algo wheel} pour qu'un courtier ne puisse pas la manipuler~?
\end{interviewq}

\begin{interviewq}[$\star\star$ \iqroles{chercheur}]\label{iq:mx:the-buy-side-trading-desk-and-its-brokers:3}
Combien de temps une rotation doit-elle tourner pour classer des courtiers écartés de 3 points de base, et qu'est-ce qui la raccourcit~?
\end{interviewq}

\begin{interviewq}[$\star\star$ \iqroles{développeur}]\label{iq:mx:the-buy-side-trading-desk-and-its-brokers:4}
Que doivent enregistrer l'OMS et l'EMS pour que la TCA et la rotation du desk fonctionnent~?
\end{interviewq}

\begin{interviewq}[$\star\star$ \iqroles{ingénieur ML}]\label{iq:mx:the-buy-side-trading-desk-and-its-brokers:5}
Feriez-vous tourner la rotation comme un bandit~? Quels sont les risques~?
\end{interviewq}

\begin{interviewq}[$\star\star\star$ \iqroles{banque}]\label{iq:mx:the-buy-side-trading-desk-and-its-brokers:6}
Vous dirigez l'activité d'algorithmes d'un courtier. La rotation d'un client vous classe cinquième sur six après un trimestre. Que faites-vous~?
\end{interviewq}
